\begin{tikzpicture}[
  x=1mm,y=1mm,
  every node/.style={font=\figurefont\small,text=BookInk},
  token/.style={book node,draw=BookRule,fill=white,minimum height=8mm,inner sep=2pt},
  span a/.style={draw=BookBlue,fill=FigureInput!10,line width=0.9pt,minimum height=7mm,inner sep=1.5pt},
  span b/.style={draw=BookAmber,fill=BookAmber!10,line width=0.9pt,minimum height=7mm,inner sep=1.5pt}
]
  \node[font=\figurefont\bfseries] at (36,28) {对象错配造成伪冲突};
  \node[token,minimum width=14mm] (ll1) at (9,13) {林舟};
  \node[token,minimum width=9mm] (join1) at (24,13) {在};
  \node[token,minimum width=13mm] (sh1) at (37,13) {上海};
  \node[token,minimum width=9mm] (verb1) at (49,13) {加入};
  \node[token,minimum width=22mm] (co1) at (66,13) {海岬科技};
  \node[span a,minimum width=20mm] (a1) at (9,-2) {来源A：人物};
  \node[span b,minimum width=26mm] (b1) at (66,-16) {来源B：组织};
  \draw[BookBlue,line width=0.9pt] (ll1.south) -- (a1.north);
  \draw[BookAmber,line width=0.9pt] (co1.south) |- (b1.north);
  \draw[book arrow,draw=BookAmber,dashed] (a1.south east) -- node[below,font=\figurefont\scriptsize,text=BookAmber] {错误对应} (b1.north west);
  \node[font=\figurefont\scriptsize,text=BookMuted,align=center] at (36,-29) {标签落在不同实体上\\不能作为同一对象的类别冲突};

  \draw[BookRule,line width=0.7pt] (80,-35) -- (80,32);

  \node[font=\figurefont\bfseries] at (121,28) {对齐后识别真实分歧};
  \node[token,minimum width=14mm] (ll2) at (94,13) {林舟};
  \node[token,minimum width=9mm] (join2) at (109,13) {在};
  \node[token,minimum width=13mm] (sh2) at (122,13) {上海};
  \node[token,minimum width=9mm] (verb2) at (134,13) {加入};
  \node[token,minimum width=22mm] (co2) at (151,13) {海岬科技};
  \node[span a,minimum width=20mm] (a2) at (94,-2) {来源A：人物};
  \node[span b,minimum width=20mm] (b2) at (94,-16) {来源B：组织};
  \draw[BookBlue,line width=0.9pt] (ll2.south) -- (a2.north);
  \draw[BookAmber,line width=0.9pt] (a2.south) -- (b2.north);
  \draw[BookTeal,line width=1pt,-{Stealth[length=2mm]}] (ll2.south east) to[out=-25,in=-155] node[below,font=\figurefont\scriptsize,text=BookTeal] {加入：主体$\to$组织} (co2.south west);
  \node[font=\figurefont\scriptsize,text=BookMuted,align=center] at (128,-29) {同一跨度上的“人物/组织”\\才是需要融合的类别分歧};
\end{tikzpicture}
